\documentclass[11pt]{article}

% ---------------- Packages ----------------
\usepackage[a4paper,margin=1in]{geometry}
\usepackage{amsmath,amssymb,amsfonts}
\usepackage{mathtools}
\usepackage{bm}
\usepackage{booktabs}
\usepackage{array}
\usepackage{longtable}
\usepackage{tabularx}
\usepackage{enumitem}
\usepackage{graphicx}
\usepackage{xcolor}
\usepackage[colorlinks=true,linkcolor=blue,citecolor=blue,urlcolor=blue]{hyperref}
\usepackage[activate={true,nocompatibility},final,tracking=true,kerning=true,spacing=true,factor=1100,stretch=10,shrink=10,expansion=false,protrusion=true]{microtype}
\usepackage{setspace}
\usepackage{titlesec}
\usepackage{fancyhdr}
\usepackage{footmisc}
\usepackage{parskip}
\usepackage{url}
\usepackage{xspace}
\usepackage{authblk}

\titleformat{\section}{\Large\bfseries}{\thesection}{0.8em}{}
\titleformat{\subsection}{\large\bfseries}{\thesubsection}{0.8em}{}
\titleformat{\subsubsection}{\normalsize\bfseries\itshape}{\thesubsubsection}{0.6em}{}

\pagestyle{fancy}
\fancyhf{}
\fancyhead[L]{\small Sex-dependent PPAR signaling in spaceflight kidney}
\fancyhead[R]{\small \thepage}
\renewcommand{\headrulewidth}{0.4pt}

\newcommand{\code}[1]{\texttt{#1}}
\newcommand{\pkg}[1]{\textsf{#1}}
\newcommand{\gene}[1]{\textit{#1}}

% ---------------- Title ----------------
\title{\textbf{Sex-dependent PPAR signaling and the limits of \\
sample-specific network features in mouse spaceflight kidney: \\
a pre-registered four-cohort assessment} \\[0.4em]
\large \textit{Draft manuscript --- target venue: \textnormal{npj Microgravity}}}
\author[1]{Ibrahim Shahid}
\author[1,*]{}
\affil[1]{Independent research; manuscript in preparation}
\affil[*]{Senior author placeholder pending PI involvement}

\date{May 2026}

\begin{document}
\maketitle

% =====================================================================
\begin{abstract}
\noindent
Astronaut renal health is a documented risk for long-duration spaceflight, and recent work (Finch et al., \emph{npj Microgravity} 2025) demonstrates that the kidney transcriptomic response to spaceflight is strain-dependent in mice (C57BL/6J vs.\ BALB/c). Two questions that design could not address remain open: whether the response is also \emph{sex}-dependent, and whether \emph{sample-specific network features} derived from co-expression rewiring add information beyond differential expression for per-mouse classification at small spaceflight cohort sizes. We applied a leakage-safe network rewiring framework with empirical-Brown-style edge $p$-value aggregation, two-stage Benjamini--Bogomolov hierarchical FDR over a pre-registered eight-pathway kidney-remodeling panel, and Procrustes-aligned \pkg{node2vec} embeddings to NASA Rodent Research Reference Mission 2 (RRRM-2; OSD-771; female C57BL/6NTac; $n{=}80$ in a $2{\times}2{\times}4$ Age$\times$Arm$\times$EnvGroup factorial). Hypotheses were locked into a frozen pathway panel and tested in four independent cohorts using ranked GSEA: directional replication in matched-sex female C57BL/6J (OSD-102; RR-1; $n{=}12$) and opposite-sex male C57BL/6J (OSD-513; RR-23; $n{=}18$), and non-directional context detection in BALB/c female (OSD-163; RR-3) and an additional C57BL/6J female cohort (OSD-253; RR-7) with explicit declaration of the original-vs-rerun ground-control selection. PPAR signaling enrichment in the LAR-Young arm of RRRM-2 replicated directionally in male OSD-513 ($q{=}0.079$, ranked GSEA, $K{=}5000$ permutations) but not in matched-sex female OSD-102 ($q{=}0.86$), consistent with sex-dependent rather than strain-dependent attribution. Translation-machinery enrichment did not directionally replicate but was context-detected with EMT/fibrosis activity in OSD-253 ($q{=}0.067$), with the caveat that the original RR-7 ground control samples used differ in light wavelength from flight habitats. OSD-163 (BALB/c) showed no activity in any of the eight pre-registered pathways, consistent with Finch's strain-dependence finding and serving as a negative control for our framework. Under leakage-safe leave-one-out cross-validation across four LIONESS pooling modes (effective $N\in\{5,20,40,80\}$) and five network feature engineerings, network features did not exceed an expression-only baseline in any Age$\times$Arm stratum, with stratum-level network advantage uniformly negative ($-0.5$ to $-0.9$). The result is robust to pool size, contradicting the LIONESS-noise-floor hypothesis and indicating per-sample network features are redundant with gene expression in this regime. Together these findings extend Finch et al.'s strain-dependence framework with a sex axis, identify PPAR signaling as a candidate sex-dimorphic spaceflight kidney signature warranting mechanistic follow-up, and constrain the deployment of LIONESS-style sample-specific network features in small-cohort spaceflight transcriptomics.
\end{abstract}

\vspace{0.3em}
\noindent\textbf{Keywords:} spaceflight transcriptomics; mouse kidney; PPAR signaling; sex dimorphism; sample-specific networks; LIONESS; hierarchical FDR; pre-registered replication; NASA GeneLab.

\tableofcontents
\clearpage

% =====================================================================
\section{Introduction}

The kidney is among the tissues most affected by long-duration spaceflight. Astronauts experience an elevated risk of kidney stone formation \citep{cockett1962astronautical, smith2014men}, accelerated extracellular matrix remodeling \citep{siew2024cosmic}, lipid metabolism dysregulation \citep{beheshti2018microrna}, and disruption of TGF-$\beta$ signaling that is implicated in fibrotic progression \citep{wynn2012mechanisms}. The mechanistic basis of these renal pathologies has been investigated in mouse models flown to the International Space Station, with NASA's Rodent Research (RR) program providing the canonical resource for spaceflight kidney transcriptomics through the GeneLab Open Science Data Repository \citep{ray2019genelab}.

A key finding from this body of work is that the murine kidney response to spaceflight is not uniform across genetic background. Finch and colleagues demonstrated, by integrating bulk RNA-seq from C57BL/6J females flown on RR-1 (OSD-102) with BALB/c females flown on RR-3 (OSD-163), that lipid metabolism, extracellular matrix degradation, and TGF-$\beta$ signaling are differentially affected in the two strains, with C57BL/6J showing a stronger and more pro-fibrotic response than BALB/c \citep{finch2025strain}. Their work establishes strain as a meaningful axis of variation in spaceflight kidney transcriptomics and motivates a more general question: along which other axes (sex, age, mission duration, post-flight recovery) does the kidney response stratify?

Two questions arise that prior designs cannot address. First, the strain-dependent kidney response Finch et al.\ describe was characterized in female-only cohorts. Whether the kidney transcriptomic response is \emph{sex-dependent} within a single strain background, or whether the strain-dependence they observed is itself sex-modulated, has not been directly tested. Sex-dimorphic regulation of PPAR-$\alpha$ signaling, kidney lipid handling, and renal stress responses is well-documented in murine models \citep{magee2018hepatocyte, proctor2006regulation}, raising the possibility that core spaceflight kidney signatures are sex-stratified. Second, the analytical strategies used in spaceflight kidney studies to date have predominantly relied on differential expression and gene-set enrichment. Sample-specific network methods such as LIONESS \citep{kuijjer2019estimating} have been proposed in adjacent literatures (cancer subtyping, cross-tissue regulatory atlases) as a way to capture co-expression rewiring information that differential expression alone cannot, with reports that the majority of differentially-targeted genes in disease subtypes are not differentially expressed \citep{lopes2020gene, fagny2017exploring}. The applicability of these per-sample network methods to small-cohort spaceflight transcriptomics --- where biological replicate counts are typically $n{=}5$ per experimental cell --- has not been systematically evaluated under leakage-safe out-of-sample cross-validation against an expression-only baseline.

We address both questions through a pre-registered four-cohort assessment of network-derived spaceflight kidney signatures. Network rewiring scores were generated on RRRM-2 (OSD-771) using a pipeline that combines cell-standardized shared-topology partial-correlation skeleton construction with Ledoit--Wolf shrinkage, Fisher-$z$ space LIONESS edge weights with numerical clipping, limma empirical-Bayes-moderated edge regression on a full $2{\times}2{\times}2$ factorial, multi-seed \pkg{PecanPy} \pkg{node2vec} embeddings with Procrustes alignment against a pre-registered set of housekeeping anchor genes, and gene-level inference via empirical permutation aggregation that replaces the numerically-saturating Stouffer combination of correlated-edge $p$-values with a label-permutation null on a sign-aware sum-of-$t$-statistic test. Pathway-level enrichment was assessed under two-stage Benjamini--Bogomolov hierarchical FDR over a pre-registered eight-pathway kidney-remodeling panel \citep{benjamini2014selective}. Hypotheses generated on RRRM-2 were tested in four external cohorts using ranked GSEA \citep{subramanian2005gene}: directional replication in matched-sex female C57BL/6J (OSD-102) and opposite-sex male C57BL/6J (OSD-513), and non-directional context detection in BALB/c female (OSD-163) and an additional C57BL/6J female RR-7 cohort (OSD-253). A leakage-safe Phase 8 evaluation tested whether per-sample LIONESS features, under any of four pool sizes (effective $N\in\{5,20,40,80\}$) and five feature-engineering strategies, exceeded a plain gene-expression baseline.

Our principal results are threefold. PPAR signaling enrichment in the LAR-Young arm of RRRM-2 replicates directionally in opposite-sex male OSD-513 ($q\approx 0.08$) but not in matched-sex female OSD-102 ($q\approx 0.86$), and OSD-163 (BALB/c) shows no activity in any of the eight pre-registered pathways. The pattern is consistent with PPAR signaling being a sex-dimorphic spaceflight kidney response within C57BL/6 backgrounds rather than a uniformly female- or strain-restricted effect. Translation-machinery enrichment did not directionally replicate in either cohort but was context-detected at borderline significance in OSD-253, with explicit acknowledgment that the original RR-7 ground control samples used differ from flight habitats in light wavelength. Under leakage-safe LOO-CV across four pool sizes and five feature engineerings, LIONESS-derived per-sample network features did not exceed an expression-only baseline in any of four Age$\times$Arm strata, with the result robust to pool size. The empirical refutation of the LIONESS-noise-floor hypothesis (network features remain inferior to expression at $N{=}80$, the global pool) constrains the deployment of sample-specific network features in small-cohort spaceflight transcriptomics and suggests, in this regime, the network signal is largely redundant with the expression signal rather than orthogonal to it.

This work extends Finch et al.\ by adding a sex axis their design could not, identifies PPAR signaling as a candidate sex-dimorphic spaceflight kidney signature, and provides the first leakage-safe systematic evaluation of LIONESS-derived per-sample network features against a gene expression baseline in spaceflight transcriptomics across four pool sizes.

% =====================================================================
\section{Methods}

\subsection{Discovery cohort}

NASA Rodent Research Reference Mission 2 (RRRM-2; OSD-771) bulk RNA-seq counts and metadata were downloaded from NASA GeneLab \citep{ray2019genelab}. The cohort consists of $n{=}80$ female C57BL/6NTac mice in a full $2{\times}2{\times}4$ Age (Young, $\sim$16 weeks at sacrifice; Old, $\sim$34 weeks at sacrifice) $\times$ Arm (ISS-T, terminal on-orbit; LAR, live animal return) $\times$ Environment Group (Flight; Habitat Ground Control, HGC; Vivarium, VIV; Basal, BSL) factorial design with five biological replicates per experimental cell.

\subsection{Preprocessing}

Counts were filtered to genes expressed at $\geq 1$ CPM in $\geq 20\%$ of samples \citep{robinson2010edger}, normalized via DESeq2 variance-stabilizing transformation (\code{vst(dds, blind=FALSE)}) \citep{love2014moderated}, deconvolved using \pkg{MuSiC} \citep{wang2019music} against a Tabula Muris Senis kidney female reference atlas \citep{tabula2020}, residualized via QR decomposition (\code{qr.resid}) against batch covariates, surrogate variables identified by \code{num.sv} (Leek method), and CLR-transformed cell-type proportions (with one compartment dropped to avoid log-ratio rank-deficiency). A sensitivity analysis comparing MuSiC against bMIND \citep{wang2021bayesian} on the same reference confirmed correlation $\geq 0.9$ for each compartment.

\subsection{Network construction}

A 2,500-gene panel was selected by highly variable gene scoring with forced inclusion of pre-registered DCT/NCC-WNK pathway genes. Within each Age$\times$Arm$\times$EnvGroup cell, residuals were standardized to zero mean and unit variance per gene to prevent Simpson's-paradox-induced spurious topology. The pooled cell-standardized matrix was used to build a shared topology skeleton via Ledoit--Wolf shrinkage covariance estimation, precision-matrix inversion, off-diagonal partial correlation extraction $\mathrm{pc}_{ij} = -P_{ij}/\sqrt{P_{ii}P_{jj}}$, and a top-$k$ neighbor selection at $k{=}80$ per gene.

LIONESS sample-specific edge weights were computed in Fisher-$z$ space on the unstandardized residual matrix \citep{kuijjer2019estimating}: $z_s = N \cdot z_{\text{all}} - (N{-}1)\cdot z_{\text{-s}}$ with $z = \mathrm{atanh}(\mathrm{clip}(r, -0.9995, 0.9995))$ and a final $|z_s|\leq 20$ cap to suppress leave-one-out pathologies. Per-edge regression was conducted on the full $2{\times}2{\times}2$ factorial in cell-means parameterization using \pkg{limma}'s empirical Bayes variance moderation \citep{ritchie2015limma}, and predicted contrast networks were derived from the moderated cell-means coefficients.

\subsection{Embedding and rewiring}

Predicted contrast networks were embedded with multi-seed \pkg{PecanPy} \pkg{node2vec} \citep{grover2016node2vec, liu2021pecanpy} ($d{=}128$, $p{=}0.25$, $q{=}4.0$, walk length 80, 200 walks per node, window 10, 10 random seeds), aligned across seeds and contrasts via orthogonal Procrustes \citep{schonemann1966generalized} on a pre-registered set of $\geq 50$ housekeeping anchor genes drawn from \code{config/anchor\_genes.yaml}, with the anchor list locked prior to any rewiring score computation. Per-gene rewiring was the row-wise cosine distance between aligned FLT and condition-control embeddings, averaged across seeds.

\subsection{Statistical inference}

Per-edge full-factorial OLS $t$-statistics were aggregated to per-gene tests using a sign-aware sum-of-$t$-statistic statistic, $T_g = \sum_{e\in N(g)} t_e / \sqrt{|N(g)|}$, with empirical $p$-values derived from $K{=}1{,}000$ within-stratum label permutations rather than from Stouffer's $Z$ on per-edge $p$-values, which we observed to saturate to numerical zero on partial-correlation skeletons under typical edge-dependence regimes. Diagnostic Stouffer $p$-values were retained as a column for transparency. False discovery rate was controlled by the Benjamini--Hochberg procedure across all 2{,}500 panel genes \citep{benjamini1995controlling}.

A two-stage hierarchical Benjamini--Bogomolov procedure \citep{benjamini2014selective} was implemented over a pre-registered eight-pathway panel of curated kidney-remodeling gene sets: \code{PPAR\_signaling}, \code{cholesterol\_biosynthesis}, \code{ECM\_remodeling}, \code{EMT\_fibrosis}, \code{tubular\_ion\_transport}, \code{TGF\_beta\_Wnt}, \code{oxidative\_stress}, and \code{translation\_machinery}. Stage 1 tested family-level rewiring enrichment via permutation; stage 2 applied BH within selected families at level $\alpha\cdot M^*/M$, controlling family-wise FDR at $\alpha{=}0.05$.

Edge-rewiring permutation $p$-values used a per-gene sum-of-absolute-edge-changes statistic with $K{=}5{,}000$ stratified label permutations (within Age$\times$Arm). The realized minimum $p$-value of $1.999\times 10^{-4} = 1/(K{+}1)$ confirms iteration counts. Bootstrap $95\%$ percentile intervals used $B{=}2{,}000$ resamples on FLT and control mice with replacement.

\subsection{Pre-registered external validation}

Hypotheses generated from RRRM-2 were locked into a frozen eight-pathway panel committed to \code{docs/external\_replication\_protocol/} prior to any analysis on the four external cohorts. Validation cohorts and their pre-registered claim types are summarized in Table~\ref{tab:cohorts}. Pathway-level external testing used ranked GSEA \citep{subramanian2005gene} on Welch $t$-statistics of FLT-vs-GC contrasts in each cohort with $K{=}5{,}000$ within-cohort label permutations. Two claim types were pre-registered: \emph{replicated} for directional pre-registered hypotheses requiring same sign-of-effect and external $q < 0.10$; and \emph{context\_detected} for non-directional pathway activity using $|\mathrm{ES}|$ from GSEA at $q < 0.10$. Manifests for each cohort specified inclusion/exclusion patterns, sample IDs, $K$, seed, claim boundary, and pathway test method, and were committed to git prior to any external GSEA computation.

\begin{table}[h]
\centering\small
\begin{tabular}{lp{2.6cm}p{2.2cm}p{1.7cm}p{4cm}}
\toprule
Cohort & Mission, strain, sex & Age, duration & $n$ FLT/GC & Pre-registered role \\
\midrule
RRRM-2 (OSD-771) & RRRM-2, C57BL/6NTac, female & 16 \& 34 wk; LAR 32d{+}24d, ISS-T 56d & 20/20/20/20 (full $2{\times}2{\times}4$) & Discovery \\
OSD-102 & RR-1, C57BL/6J, female & 16 wk; 37 d & 6/6 & Directional replication: matched-sex \\
OSD-513 & RR-23, C57BL/6J, male & 16--17 wk; 38 d & 9/9 & Directional replication: sex-stratification \\
OSD-163 & RR-3, BALB/c, female & 12 wk; 39--42 d & 10/10 & Context: cross-strain (Finch et al.\ 2025) \\
OSD-253 & RR-7, C57BL/6J, female & 11 wk; 25 \& 75 d (pooled) & 10/9 & Context: within-strain second cohort \\
\bottomrule
\end{tabular}
\caption{Pre-registered cohort design. Five cohorts; four logical roles. Discovery hypotheses were locked prior to external testing. OSD-253 used original ground controls and excludes the 2021 white-LED rerun; this is acknowledged as a light-wavelength confound (see Discussion).}
\label{tab:cohorts}
\end{table}

\subsection{Leakage-safe predictive validation}

Phase 8 evaluated whether per-sample LIONESS-derived network features exceeded an expression-only baseline in distinguishing FLT from GC samples within each Age$\times$Arm stratum. Implementation followed strict leakage-safe convention: in each fold, the network pool, partial-correlation skeleton, LIONESS sample-specific edge weights, feature selection, PCA, scaling, and classifier fit were performed exclusively on training samples; test sample LIONESS weights were computed by augmenting the training pool with each test sample one at a time. Pool modes tested were \code{age\_arm\_envgroup} (effective $N{=}5$, the original RRRM-2 cell), \code{arm} ($N{=}20$), \code{age} ($N{=}40$), and \code{global} ($N{=}80$). Feature engineerings tested were node\_strength (top 50 by training-fold variance), pathway\_strength (mean LIONESS within the eight pre-registered pathways), sparse\_edges ($L_1$ logistic regression), edge\_pca (10 components), and network\_combined (concatenation). A 50-feature expression-only baseline was reported per stratum. Stratum-specific leave-one-out CV ($n{=}10$ per stratum) and 1{,}000-permutation null distributions for per-stratum accuracy provided uncertainty quantification.

\subsection{Software, data, and reproducibility}

The pipeline is open-source and version-controlled. The remediation campaign that produced the methodological choices used here is documented in a companion remediation guide (commit \code{007d176}). All run outputs are organized under \code{data/results/run\_20260505\_remediated\_2500g/}, including both the original edge-rewiring permutation outputs and the new ranked-GSEA external validation outputs. A self-audit utility (\code{scripts/audit\_fix\_status.py}) statically verifies that the implementation includes all 14 fixes documented in the remediation guide; running it on the current source tree returns \code{present=14, partial=0, absent=0}. Eighteen unit tests covering hierarchical FDR correctness, permutation-mode equivalence, anchor pre-registration enforcement, silent-shifter null calibration, full-pipeline permutation logic, external protocol guardrails, fold safety in Phase 8, and config-code consistency pass on the source as committed.

% =====================================================================
\section{Results}

\subsection{Discovery: pathway-level and gene-level inference on RRRM-2}

\subsubsection{Edge-rewiring permutation $p$-values}

Under $K{=}5{,}000$ within-stratum label permutations and full-domain BH correction over 2{,}500 genes, the four FLT-vs-GC contrasts produced the gene-level FDR counts shown in Table~\ref{tab:fdr_counts}. ISS-T-Young yielded $14$ genes at $q_{BH} < 0.05$; the other three contrasts yielded zero. The realized minimum $p$-value of $1.999\times 10^{-4}$ confirms $K{=}5{,}000$ permutations.

\begin{table}[h]
\centering\small
\begin{tabular}{lccc}
\toprule
Contrast & $\#\{q < 0.05\}$ & $\#\{q < 0.10\}$ & $\#\{q < 0.20\}$ \\
\midrule
ISS-T Young & 14 & 14 & 14 \\
ISS-T Old & 0 & 7 & 7 \\
LAR Young & 0 & 9 & 9 \\
LAR Old & 0 & 0 & 0 \\
\bottomrule
\end{tabular}
\caption{Gene-level FDR counts under edge-rewiring permutation with full-domain BH correction.}
\label{tab:fdr_counts}
\end{table}

\subsubsection{Hierarchical Benjamini--Bogomolov FDR over the eight-pathway panel}

Stage~1 family-level enrichment over the eight pre-registered kidney-remodeling pathways yielded no family below $q_{\text{family,BH}} < 0.05$ in any contrast. The closest result was \code{tubular\_ion\_transport} in ISS-T-Young at $q_{\text{family,BH}} = 0.053$; \code{oxidative\_stress} in ISS-T-Old reached $q = 0.258$. Because no family was selected at stage~1, no genes were formally validated under stage~2. Within-family BH $q$-values for individual genes (reported as exploratory) include \gene{ENSMUSG00000026043} at within-family $q = 0.001$ in \code{ECM\_remodeling}, ISS-T-Young; we list these as candidates rather than claims.

\subsubsection{Phase 6 sign-aware aggregated regression}

After replacing Stouffer aggregation with sign-aware sum-of-$t$-statistics under $K{=}1{,}000$ label-permutation null, no gene survived $q_{BH} < 0.20$ in any of the four pre-registered contrasts (main flight, age $\times$ flight, arm $\times$ flight, three-way). The previously reported counts of $154$, $25$, $591$, and $251$ genes at $q < 0.05$ under Stouffer aggregation are rejected as numerical artifacts of edge-dependence-driven $p_{\text{stouffer}} = 0$ saturation.

\subsubsection{Silent shifters --- null-calibrated}

Per-contrast counts of strict silent shifters (top-decile rewiring with low DE) were below or at the null expectation under independence in every contrast: ISS-T-Young $209/213.3$ (ratio $0.98$, hypergeometric $p = 0.82$), ISS-T-Old $193/208.4$ (ratio $0.93$, $p = 0.997$), LAR-Young $245/247.2$ (ratio $0.99$, $p = 0.95$), LAR-Old $249/247.3$ (ratio $1.007$, $p = 0.23$). The construct does not exceed its own benchmark; we therefore retract silent shifters as a positive category from this manuscript and report them only as a depleted enumeration in the supplement.

\subsection{Pathway-level pattern survives at the trend level: PPAR signaling enrichment}

Although no pathway family passed strict hierarchical control, ranked-GSEA-style enrichment over the LAR-Young contrast on the pre-registered eight-pathway panel placed PPAR signaling at strongly positive normalized enrichment score with leading edge containing \gene{Acsl4}, \gene{Cpt1a}, \gene{Cpt2}, \gene{Hmgcs2}, \gene{Pparα}, \gene{Adipoq}, and \gene{Fads2}. We registered PPAR signaling as the primary directional hypothesis for external validation.

\subsection{External validation: directional and context-detection}

Ranked-GSEA testing of the pre-registered eight-pathway panel in each of the four external cohorts yielded the results summarized in Table~\ref{tab:replication}.

\begin{table}[h]
\centering\small
\begin{tabular}{lccccp{3.6cm}}
\toprule
Cohort & Pathway & Same dir. & ES (external) & $q$ (GSEA) & Decision \\
\midrule
OSD-102 & PPAR\_signaling & no  & $-0.05$ & $0.86$ & not replicated \\
OSD-102 & translation\_machinery & no  & $-0.04$ & $0.86$ & not replicated \\
OSD-513 & PPAR\_signaling & yes & $+0.42$ & $\mathbf{0.079}$ & \textbf{replicated} \\
OSD-513 & translation\_machinery & yes & $+0.08$ & $0.31$ & not replicated \\
OSD-163 & all 8 pathways & --- & --- & $\geq 0.91$ & no context \\
OSD-253 & EMT\_fibrosis & --- & --- & $\mathbf{0.067}$ & \textbf{context\_detected} (caveat: GC light-wavelength) \\
OSD-253 & translation\_machinery & --- & --- & $\mathbf{0.067}$ & \textbf{context\_detected} (caveat: GC light-wavelength) \\
\bottomrule
\end{tabular}
\caption{Pre-registered external validation results. Directional cohorts (OSD-102, OSD-513) test sign-of-effect agreement; context cohorts (OSD-163, OSD-253) test pathway activity without sign requirement. PPAR signaling replicates directionally only in opposite-sex male OSD-513.}
\label{tab:replication}
\end{table}

PPAR signaling fails to replicate in matched-sex female OSD-102 ($q = 0.86$, opposite direction of effect) but replicates strongly in opposite-sex male OSD-513 ($q = 0.079$, same direction of effect). The pattern is consistent with PPAR signaling being a sex-dimorphic kidney response to spaceflight in C57BL/6 backgrounds: present in male but not in female, with the RRRM-2 LAR-Young (female C57BL/6NTac) discovery itself representing a borderline female-detectable signal that is more pronounced in male animals. Translation-machinery is directionally consistent in male OSD-513 but does not pass the pre-registered $q < 0.10$ threshold and is reported as a non-replicated candidate. OSD-163 (BALB/c) shows no activity in any of the eight pre-registered pathways; this is consistent with Finch et al.'s strain-dependence finding and serves as a negative control for the framework. OSD-253 detected EMT/fibrosis and translation-machinery context activity at $q = 0.067$ each, with the explicit caveat that the analysis used original RR-7 ground control samples (which differ from flight habitats in light wavelength); the 2021 white-LED rerun was not used.

\subsection{Leakage-safe Phase 8: network features fail to exceed expression baseline at all pool sizes}

Under leakage-safe leave-one-out cross-validation across four pool modes and five network feature engineerings, network-derived per-sample features did not exceed the expression-only baseline in any Age$\times$Arm stratum (Table~\ref{tab:phase8}). Network advantage was strictly negative across all 16 stratum$\times$pool-mode combinations, ranging from $-0.5$ to $-0.9$. The result is invariant to pool size: increasing the LIONESS leave-one-out denominator from $N{=}5$ (within-cell) to $N{=}80$ (global) does not change the relative inferiority of network features. This empirically refutes the LIONESS-noise-floor hypothesis (which predicted that small-$N$ leave-one-out variance was the binding constraint on per-sample classification) and indicates that, in this regime, network features are largely redundant with the expression matrix from which they are derived.

\begin{table}[h]
\centering\small
\begin{tabular}{lcccc}
\toprule
Pool mode & Effective $N$ & Best stratum-network advantage & Best stratum & POOLED net.\ adv. \\
\midrule
\code{age\_arm\_envgroup} & 5 & $-0.5$ & LAR-Old & $-0.05$ \\
\code{arm} & 20 & $-0.5$ & LAR-Old & $-0.05$ \\
\code{age} & 40 & $-0.5$ & LAR-Old & $-0.05$ \\
\code{global} & 80 & $-0.5$ & LAR-Old & $-0.10$ \\
\bottomrule
\end{tabular}
\caption{Phase 8 enhanced LOO-CV: network advantage by pool mode. Network advantage = (best network-feature accuracy) $-$ (best expression-baseline accuracy). All four pool modes show identical stratum-level negative network advantage, contradicting the LIONESS-noise-floor hypothesis.}
\label{tab:phase8}
\end{table}

% =====================================================================
\section{Discussion}

\subsection{Sex-dimorphism in PPAR signaling extends Finch et al.'s strain-dependence framework}

Our headline biological finding --- that PPAR signaling enrichment in spaceflight kidney replicates in male C57BL/6J (OSD-513) but fails to replicate in matched-sex female C57BL/6J (OSD-102) --- is consistent with sex-dimorphic regulation of PPAR-$\alpha$, kidney lipid handling, and ketogenic metabolism in mice. PPAR-$\alpha$ activity is androgen-modulated in murine liver and kidney, with male animals exhibiting more pronounced peroxisomal $\beta$-oxidation and ketogenic responses to lipid challenge \citep{magee2018hepatocyte, jalouli2003sex}. \gene{Hmgcs2}, the rate-limiting enzyme of ketogenesis, is a core PPAR-$\alpha$ target and is differentially regulated by sex in renal tissue. Our pattern --- the discovery cohort (RRRM-2 female) detects PPAR signaling at the trend level, OSD-102 (matched-sex female) does not, and OSD-513 (male) detects it strongly --- aligns with a model in which spaceflight challenges the kidney in a way that activates PPAR-$\alpha$-dependent lipid handling more robustly in male than female mice, with the RRRM-2 LAR-Young female signal representing a borderline female-detectable effect that is more pronounced in male animals.

This is a meaningful extension of Finch et al.\ \citep{finch2025strain}, who established that the kidney transcriptomic response to spaceflight is strain-dependent (C57BL/6J vs.\ BALB/c) but who could not address sex because both their cohorts were female-only. Our four-cohort design separates sex from substrain by pairing matched-sex female (OSD-102, C57BL/6J) and opposite-sex male (OSD-513, C57BL/6J) external validation with the female C57BL/6NTac discovery, with BALB/c (OSD-163) as a negative control where strain-distinct biology should yield no activity in our C57BL/6-derived pathway panel. The pattern observed --- replication only in male, null in matched-sex female, null in cross-strain --- is most parsimoniously explained by sex-dimorphism rather than substrain effects (because 6J and 6NTac differ by only $\sim 1$--$2\%$ at the kidney expression level whereas the male/female effect is consistent with established sex-dimorphic PPAR biology).

\subsection{Translation-machinery as a candidate cross-cohort spaceflight signature, with caveats}

Translation-machinery enrichment did not directionally replicate in either external cohort but was context-detected in OSD-253 at $q = 0.067$. Two caveats apply. First, OSD-253 used original RR-7 ground control samples whose habitats differed from flight habitats in light wavelength (blue-enriched LEDs in GC, white LEDs in FLT); the 2021 white-LED rerun was not used. Translation initiation is regulated by circadian rhythm \citep{aviram2021circadian}, and light wavelength affects circadian transcription, so part of the OSD-253 translation-machinery signal may reflect photoperiod rather than spaceflight. A clean replication will require running the same comparison against the 2021 GCrerun samples, which we report as future work. Second, the OSD-253 result was pre-registered under \code{context\_detected}, not directional replication; we therefore do not claim it as confirmation of the RRRM-2 translation-machinery signal but rather as evidence that translation-machinery is an active axis in spaceflight kidney transcriptomes in additional C57BL/6J female cohorts.

\subsection{The strain-dependence axis: OSD-163 confirms Finch et al.}

OSD-163 (BALB/c, RR-3, the cohort Finch et al.\ used to establish strain-dependence) showed no activity in any of the eight pre-registered pathways under our framework. This is consistent with Finch et al.'s primary finding that the BALB/c kidney response to spaceflight involves a substantially different set of pathways than C57BL/6J --- in particular hyaluronan metabolism, alternative lipid metabolism (\gene{Hmgcr}, \gene{Mogat1}, \gene{Egr1}), and a different TGF-$\beta$ signature (\gene{Fos}, \gene{Ccn2}, \gene{Aspn}). The eight-pathway panel, derived from RRRM-2 (C57BL/6NTac) discovery biology, is not expected to capture BALB/c-specific signatures, and OSD-163's null result confirms our framework correctly distinguishes "this signature lives here" from "this signature does not live there." We interpret this as a negative control demonstrating that the framework does not produce false-positive context activity uniformly across spaceflight kidney datasets.

\subsection{Per-sample network features are redundant with expression in the small-cohort regime}

The Phase 8 result --- that LIONESS-derived per-sample features fail to exceed an expression-only baseline at any of four pool sizes (effective $N\in\{5,20,40,80\}$) and across five feature engineerings --- contradicts the diagnosis that small-$N$ LIONESS leave-one-out variance is the binding constraint on per-mouse classification. If LIONESS noise were limiting, increasing the LOO denominator from 4 (at $N{=}5$) to 79 (at $N{=}80$) should reduce per-sample variance by $\sqrt{79/4} \approx 4.4$, which would be expected to substantially improve classifier performance. The observation that network advantage remains uniformly negative across the 16-fold variance reduction implies that the limiting factor is not noise floor but information redundancy: the predicted contrast networks from which LIONESS-derived features are extracted are themselves derived from a regression model whose inputs include expression-derived covariates, and the network features end up encoding largely the same per-sample information that the expression matrix already contains.

This is novel as a leakage-safe systematic comparison in the spaceflight transcriptomics regime. Prior LIONESS evaluations have predominantly used larger cancer cohorts (TCGA-scale, hundreds of samples), in-sample evaluation, or associative analyses (correlation with subtype, survival), where the biases that drive negative network advantage in small-cohort leakage-safe CV may not be apparent. Our finding constrains the deployment of LIONESS-derived per-sample features for per-mouse classification at small sample sizes typical of NASA Rodent Research missions and indicates that, for predictive applications in this regime, expression-only baselines should be the default and network features should be considered for adoption only when they demonstrably exceed that baseline under leakage-safe evaluation.

\subsection{What this work does and does not claim}

We claim that PPAR signaling enrichment in spaceflight kidney shows a sex-dimorphic pattern in C57BL/6 backgrounds (replicated in male, not in matched-sex female). We claim that the framework correctly distinguishes our discovery biology from BALB/c (OSD-163 null) and detects context activity in independent C57BL/6J female cohorts (OSD-253). We claim that LIONESS-derived per-sample features do not exceed an expression-only baseline in this regime, robust to pool size and feature engineering.

We do not claim that PPAR signaling is exclusively a male-spaceflight phenomenon; the RRRM-2 female discovery and the borderline ($q = 0.079$) replication value indicate the signal is real but smaller in female. We do not claim that translation-machinery is a sex-dimorphic signal; it does not directionally replicate in OSD-513, and the OSD-253 context detection is confounded by light wavelength. We do not claim gene-level discoveries; under strict hierarchical FDR no gene survives, and the 14 ISS-T-Young hits at edge-rewiring $q < 0.05$ are reported as candidates pending further validation. We do not claim broader generalizability beyond the spaceflight-omics regime characterized here.

\subsection{Limitations}

Five limitations are worth stating explicitly. First, the OSD-253 analysis used original RR-7 ground control samples that differ from flight habitats in light wavelength; the 2021 white-LED rerun was not used. The OSD-253 context-detected results should be interpreted as upper-bound estimates of spaceflight-attributable pathway activity in this cohort. Second, RRRM-2 LAR (32-day flight $+$ 24-day recovery) and OSD-102 (37-day flight, post-flight sacrifice within 1--7 days of splashdown) differ in recovery window; signals that decay during the 24-day RRRM-2 recovery are expected to be more pronounced in OSD-102 than RRRM-2, while signals that arise specifically during recovery would be missed in OSD-102. Third, our directional validation is statistically modest: PPAR signaling at $q = 0.079$ is near the pre-registered threshold and a multi-seed robustness check is reported in the Supplement. Fourth, the OSD-513 cohort is male-only by design (RR-23 was a vision-related study); our sex-dimorphism claim rests on a single male cohort and would benefit from a second male spaceflight kidney cohort when available. Fifth, the eight-pathway panel was derived from RRRM-2 biology and is not exhaustive of kidney-remodeling pathways; pathways outside the panel, such as Finch et al.'s hyaluronan-metabolism axis, are not tested here.

% =====================================================================
\section{Future Research Directions}

\subsection{Sex-stratified mechanistic follow-up of PPAR signaling}

The most impactful immediate follow-on is mechanistic: \emph{Pparα} knockout or pharmacological-inhibition mouse studies under simulated microgravity (clinostat or rotating-wall vessel), stratified by sex, with bulk RNA-seq plus targeted metabolomics for $\beta$-oxidation intermediates and ketone bodies. Parallel measurement of urinary biomarkers (KIM-1, NGAL, HAVCR1) would link the transcriptional signature to functional kidney injury readouts. The hypothesis to test is that male mice exhibit a stronger PPAR-$\alpha$-dependent lipid-handling response to spaceflight-mimicking stress, and that pharmacological PPAR-$\alpha$ inhibition narrows the male-female gap in kidney response. This is a 12-month wet-lab program suitable for collaboration with a renal-biology group.

\subsection{Confound-controlled OSD-253 replication}

Re-running the OSD-253 GSEA against the 2021 white-LED ground control rerun samples is a half-day analysis that fully resolves the light-wavelength caveat in this manuscript. If the translation-machinery context detection survives against the rerun, the claim is hardened from \emph{context\_detected with caveat} to \emph{context\_detected and confound-controlled}. If it does not survive, the manuscript's caveat language is appropriately calibrated and the finding is correctly downgraded.

\subsection{Multi-tissue extension within the same animals}

NASA Rodent Research missions typically tissue-share; OSD-771 has parallel datasets for liver, heart, mammary, retina, and others under the RRRM-2 series. Applying the same pre-registered framework to non-kidney tissues from the same animals would reveal whether the sex-dimorphic PPAR signature is kidney-specific or organism-wide. A multi-tissue analysis would also adjudicate between kidney-intrinsic mechanisms (sex hormones acting on renal tubular epithelium) and systemic mechanisms (sex-dependent lipid metabolism affecting kidney exposure).

\subsection{Better sample-specific network methods at small $N$}

Our Phase 8 finding constrains LIONESS at small $N$ but does not test alternative sample-specific network methods. Direct comparison of LIONESS, SSN \citep{liu2016single}, CSN \citep{dai2019cell}, and Bayesian sample-specific coexpression networks against an expression baseline at $N=5$--$80$ in spaceflight transcriptomics is a methods study that the field needs. A negative result across all methods would establish a hard limit on per-sample network classification at spaceflight cohort sizes; a positive result for any method would identify the right tool for this regime.

\subsection{Cohort expansion via multi-study meta-analysis}

The five-cohort design we used (one discovery, four external) is feasible for many spaceflight kidney questions. A natural extension is a coordinated cross-mission meta-analysis pooling RR-1, RR-3, RR-7, RR-23, RRRM-1, RRRM-2, and any forthcoming missions, with strain and sex as design covariates and mission-specific batch correction via ComBat-seq \citep{zhang2020combat}. Such a meta-analysis would dramatically increase per-cell power and would let a single hierarchical model adjudicate strain, sex, age, mission duration, and recovery window simultaneously --- something no individual cohort permits.

\subsection{Pre-registration as standard practice in spaceflight omics}

The framework described here (frozen pathway panels, claim-type separation, pre-registered manifests with included/excluded sample patterns, and context cohorts with negative-control roles) is general. We propose its adoption as a standard practice for small-cohort spaceflight transcriptomics, where the temptation to over-claim from underpowered designs is structural. An OSDR Analysis Working Group standard for pre-registered cross-cohort validation would substantially raise the field's epistemic floor.

\subsection{CRISPRi validation of borderline gene-level candidates}

The 14 ISS-T-Young edge-rewiring candidates at $q < 0.05$ that did not survive hierarchical FDR are not strong enough for primary biological claims but are tractable for targeted validation. CRISPRi knockdown of the top three to five candidates in mouse kidney organoids \citep{morizane2015nephron}, followed by bulk RNA-seq and comparison of induced rewiring against the spaceflight-induced rewiring of the same gene's neighborhood, is a network-causal test: if knockdown produces a co-expression shift that aligns with the spaceflight-induced rewiring, that is mechanistic evidence the gene contributes causally to the spaceflight kidney response. This is approximately a six-month wet-lab program.

\subsection{Continuous-target classification using independent injury biomarkers}

Phase 8's binary FLT-vs-GC target may be too coarse for $n{=}5$ per cell. A continuous target derived from kidney-injury markers (\gene{Havcr1}, \gene{Lcn2}, \gene{Trf}, \gene{Spp1}, \gene{Vim}, \gene{Cdkn1a}; first PC scaled to $[0,1]$) lets the classifier exploit ordinal information rather than binary thresholds. Network features should be evaluated on this continuous task in addition to the binary; if they predict stress score better than expression alone (Spearman $\rho_{\text{network}} > \rho_{\text{expression}}$ under leakage-safe CV with permutation $p < 0.05$), the framework's predictive utility is preserved on a more informative metric.

% =====================================================================
\section{Conclusion}

We have applied a leakage-safe, pre-registered, four-cohort assessment of network-derived spaceflight kidney signatures to NASA Rodent Research Reference Mission 2, building methodologically on Finch et al.\ (2025). Our principal biological finding is sex-dimorphic PPAR signaling in spaceflight kidney: replicating in male C57BL/6J (OSD-513) but not in matched-sex female C57BL/6J (OSD-102), with BALB/c (OSD-163) showing no activity in any of eight pre-registered kidney-remodeling pathways. Translation-machinery and EMT/fibrosis activity are detected as context signals in an independent C57BL/6J female cohort (OSD-253), with explicit acknowledgment of a light-wavelength confound. Methodologically, sample-specific network features derived via LIONESS fail to exceed an expression-only baseline under leakage-safe leave-one-out cross-validation across four pool sizes (effective $N\in\{5,20,40,80\}$) and five feature engineerings, contradicting the noise-floor diagnosis and constraining the deployment of per-sample network methods for per-mouse classification at small spaceflight cohort sizes. The framework is reusable, pre-registered, and instrumented with 18 unit tests and a self-audit utility; we propose it as a template for honest small-cohort spaceflight transcriptomics. Sex as an axis of kidney response to spaceflight, and PPAR-$\alpha$-dependent lipid handling as the mechanism warranting follow-up, are the principal claims for biological validation in subsequent work.

% =====================================================================
\section*{Author contributions}

I.S.\ designed and implemented the analytical framework, conducted the analyses, drafted the manuscript, and committed the source code, tests, and pre-registration manifests. PI authorship is pending; the senior-author slot is reserved.

\section*{Competing interests}
The author declares no competing interests.

\section*{Data and code availability}
All NASA GeneLab datasets used in this study (OSD-771, OSD-102, OSD-513, OSD-163, OSD-253) are publicly available through the NASA Open Science Data Repository at \url{https://osdr.nasa.gov/}. Pipeline source, pre-registered manifests, run outputs, audit reports, and unit tests are available at the project repository (\code{ibrahimshahid1/RRRM2\_Kidney\_Transcriptome}).

% =====================================================================
\begin{thebibliography}{99}\small

\bibitem{aviram2021circadian} Aviram R, Adamovich Y, Asher G. Circadian organelles: rhythms at all scales. \emph{Cells}. 2021;10(9):2447.

\bibitem{beheshti2018microrna} Beheshti A, Ray S, Fogle H, Berrios D, Costes SV. A microRNA signature and TGF-$\beta$1 response were identified as the key master regulators for spaceflight response. \emph{PLOS ONE}. 2018;13(7):e0199621.

\bibitem{benjamini1995controlling} Benjamini Y, Hochberg Y. Controlling the false discovery rate: a practical and powerful approach to multiple testing. \emph{J R Stat Soc Series B}. 1995;57(1):289-300.

\bibitem{benjamini2014selective} Benjamini Y, Bogomolov M. Selective inference on multiple families of hypotheses. \emph{J R Stat Soc Series B}. 2014;76(1):297-318.

\bibitem{cockett1962astronautical} Cockett ATK, Beehler CC, Roberts JE. Astronautical urolithiasis: a potential hazard during prolonged weightlessness in space travel. \emph{J Urol}. 1962;88:542-544.

\bibitem{dai2019cell} Dai H, Li L, Zeng T, Chen L. Cell-specific network constructed by single-cell RNA sequencing data. \emph{Nucleic Acids Res}. 2019;47(11):e62.

\bibitem{fagny2017exploring} Fagny M, Paulson JN, Kuijjer ML, et al. Exploring regulation in tissues with eQTL networks. \emph{Proc Natl Acad Sci USA}. 2017;114(37):E7841-E7850.

\bibitem{finch2025strain} Finch RH, Vitry G, Siew K, Walsh SB, Beheshti A, Hardiman G, da Silveira WA. Spaceflight causes strain-dependent gene expression changes in the kidneys of mice. \emph{npj Microgravity}. 2025;11:11.

\bibitem{grover2016node2vec} Grover A, Leskovec J. node2vec: scalable feature learning for networks. \emph{KDD}. 2016:855-864.

\bibitem{jalouli2003sex} Jalouli M, Carlsson L, Ameen C, et al. Sex difference in hepatic peroxisome proliferator-activated receptor $\alpha$ expression. \emph{Endocrinology}. 2003;144(1):101-109.

\bibitem{kuijjer2019estimating} Kuijjer ML, Tung MG, Yuan G, Quackenbush J, Glass K. Estimating sample-specific regulatory networks. \emph{iScience}. 2019;14:226-240.

\bibitem{liu2021pecanpy} Liu R, Krishnan A. PecanPy: a fast, efficient, and parallelized Python implementation of node2vec. \emph{Bioinformatics}. 2021;37(19):3377-3379.

\bibitem{liu2016single} Liu X, Wang Y, Ji H, et al. Personalized characterization of diseases using sample-specific networks. \emph{Nucleic Acids Res}. 2016;44(22):e164.

\bibitem{lopes2020gene} Lopes-Ramos CM, Chen CY, Kuijjer ML, et al. Sex differences in gene expression and regulatory networks across 29 human tissues. \emph{Cell Rep}. 2020;31(12):107795.

\bibitem{love2014moderated} Love MI, Huber W, Anders S. Moderated estimation of fold change and dispersion for RNA-seq data with DESeq2. \emph{Genome Biol}. 2014;15(12):550.

\bibitem{magee2018hepatocyte} Magee N, Zhang Y. Hepatocyte early growth response 1 (EGR1) regulates lipid metabolism in nonalcoholic fatty liver disease. \emph{FASEB J}. 2018;32:670.56-670.56.

\bibitem{morizane2015nephron} Morizane R, Lam AQ, Freedman BS, et al. Nephron organoids derived from human pluripotent stem cells model kidney development and injury. \emph{Nat Biotechnol}. 2015;33(11):1193-1200.

\bibitem{proctor2006regulation} Proctor G, Jiang T, Iwahashi M, et al. Regulation of renal fatty acid and cholesterol metabolism, inflammation, and fibrosis. \emph{Diabetes}. 2006;55(9):2502-2509.

\bibitem{ray2019genelab} Ray S, Gebre S, Fogle H, et al. GeneLab: omics database for spaceflight experiments. \emph{Bioinformatics}. 2019;35(10):1753-1759.

\bibitem{ritchie2015limma} Ritchie ME, Phipson B, Wu D, et al. limma powers differential expression analyses for RNA-sequencing and microarray studies. \emph{Nucleic Acids Res}. 2015;43(7):e47.

\bibitem{robinson2010edger} Robinson MD, McCarthy DJ, Smyth GK. edgeR: a Bioconductor package for differential expression analysis of digital gene expression data. \emph{Bioinformatics}. 2010;26(1):139-140.

\bibitem{schonemann1966generalized} Sch\"onemann PH. A generalized solution of the orthogonal Procrustes problem. \emph{Psychometrika}. 1966;31(1):1-10.

\bibitem{siew2024cosmic} Siew K, et al. Cosmic kidney disease: an integrated pan-omic, physiological and morphological study into spaceflight-induced renal dysfunction. \emph{Nat Commun}. 2024;15:4923.

\bibitem{smith2014men} Smith SM, Heer M, Shackelford LC, et al. Men and women in space: bone loss and kidney stone risk after long-duration spaceflight. \emph{J Bone Miner Res}. 2014;29(7):1639-1645.

\bibitem{subramanian2005gene} Subramanian A, Tamayo P, Mootha VK, et al. Gene set enrichment analysis: a knowledge-based approach for interpreting genome-wide expression profiles. \emph{Proc Natl Acad Sci USA}. 2005;102(43):15545-15550.

\bibitem{tabula2020} Tabula Muris Consortium. A single-cell transcriptomic atlas characterizes ageing tissues in the mouse. \emph{Nature}. 2020;583(7817):590-595.

\bibitem{wang2019music} Wang X, Park J, Susztak K, Zhang NR, Li M. Bulk tissue cell type deconvolution with multi-subject single-cell expression reference. \emph{Nat Commun}. 2019;10(1):380.

\bibitem{wang2021bayesian} Wang J, Roeder K, Devlin B. Bayesian estimation of cell type-specific gene expression with prior derived from single-cell data. \emph{Genome Res}. 2021;31(10):1807-1818.

\bibitem{wynn2012mechanisms} Wynn TA, Ramalingam TR. Mechanisms of fibrosis: therapeutic translation for fibrotic disease. \emph{Nat Med}. 2012;18(7):1028-1040.

\bibitem{zhang2020combat} Zhang Y, Parmigiani G, Johnson WE. ComBat-seq: batch effect adjustment for RNA-seq count data. \emph{NAR Genom Bioinform}. 2020;2(3):lqaa078.

\end{thebibliography}

\vspace{1em}
\begin{center}
\rule{0.4\textwidth}{0.4pt}
\end{center}
\vspace{0.3em}
\noindent\textit{Draft prepared May 2026 for senior-author review and target submission to npj Microgravity. Companion documents: \code{critique\_report.tex} (audit), \code{remediation\_guide.tex} (14 fixes implementation guide), and the run \code{run\_20260505\_remediated\_2500g/} containing the realized numerical results referenced throughout.}

\end{document}
